% START: Titel/Namen unten ersetzen, dann die Folien im Dokument bearbeiten.
% Alles hinter % ist ein Kommentar und erscheint nicht im PDF.
% Bauen: im Repository-Hauptordner `make praesentation` (siehe README.md).
\documentclass[aspectratio=169,11pt]{beamer}
% Für einen Ausdruck ohne Einblende-Schritte: oben zusätzlich `handout` setzen.
\usepackage[T1]{fontenc}
\usepackage[utf8]{inputenc}
\usepackage[ngerman]{babel}
\usepackage{lmodern}
\usepackage{amsmath,amssymb}
\usepackage{graphicx,booktabs}
\usepackage{pgfplots}
\pgfplotsset{compat=1.18}
\usepackage[authoryear,round]{natbib}
\newcommand{\quellenfussnote}[1]{\footnote[frame]{#1}}
% QUELLEN: Einheitlicher Fußnotenbefehl in Hausarbeit und Präsentation.
% Eine optionale Angabe = Fundstelle; zwei = Vorsatz (z. B. vgl.) + Fundstelle.
\usepackage{xparse}
\NewDocumentCommand{\footcite}{o o m}{%
  \quellenfussnote{%
    \IfNoValueTF{#1}{\citealp{#3}}{%
      \IfNoValueTF{#2}{\citealp[#1]{#3}}{\citealp[#1][#2]{#3}}}%
  }%
}
\usepackage{csquotes} % Deutsche Anführungszeichen mit \enquote{Text}

% DESIGN: CAU Beamer Theme von Zhuo Yue; Herkunft/Lizenz siehe CAU-QUELLE.md.
% Originaldesign mit CAU-Grün, Navigationspunkten und schlanker Fußzeile.
% Alternative Farbe z. B. CAU_SCI (Blau); die Optionen stehen in CAU.sty.
\usepackage{xkeyval}
\usepackage[color=CAU]{CAU}
\colorlet{matheblau}{use_color} % Auch die Funktionsgraphen in Themenfarbe.
% Deutsche Abschnittsübersichten; zeigen den aktuellen Abschnitt hervorgehoben.
\AtBeginSection[]{%
  \begin{frame}{Unser Weg durch die Stunde}
    \tableofcontents[currentsection,hideallsubsections]
  \end{frame}
}
\setbeamertemplate{navigation symbols}{}
\setbeamertemplate{caption}[numbered]
\setbeamertemplate{bibliography item}[text]
% Quellenfußnoten auch in Spalten unten auf der gesamten Folie ausgeben.
\setbeamerfont{footnote}{size=\scriptsize}
% Sprecherhinweise stehen im Code. Für einen eigenen Notizen-Ausdruck ersetzen
% durch \setbeameroption{show only notes}; anschließend separat kompilieren.
\setbeameroption{hide notes}

% BEARBEITEN: gemeinsamer Titel, beide vollständigen Namen, Kurs und Datum.
\title[Quadratische Funktionen]{Quadratische Funktionen}
\subtitle{Beispielvorlage für eine Mathematik-Unterrichtsstunde}
\author[Name 1 \& Name 2]{Vorname Nachname 1 \and Vorname Nachname 2}
\institute[JLU Gießen]{Justus-Liebig-Universität Gießen\\Mathematik · Kurs M12}
\date[Datum]{Abgabedatum / Datum der Unterrichtsstunde}

\begin{document}
\begin{frame}
  \titlepage
  \note{Person 1: begrüßen. Person 2: Leitfrage und Ablauf vorstellen.}
\end{frame}

\begin{frame}{Unser Weg durch die Stunde}
  \tableofcontents
  % BEARBEITEN: Abschnitte unten ändern; diese Übersicht passt sich an.
\end{frame}

\section[Einstieg]{Einstieg und Lernziele}
\begin{frame}{Leitfrage und Lernziele}
  % BEARBEITEN: konkretes Thema, Vorwissen und überprüfbare Lernziele.
  \begin{block}{Leitfrage}
    Wie finden wir die Nullstellen einer quadratischen Funktion?
  \end{block}
  Nach der Stunde könnt ihr \ldots
  \begin{itemize}
    \item einen Funktionsgraphen mit seiner Gleichung verbinden,
    \item Nullstellen durch Faktorisieren bestimmen,
    \item das Ergebnis rechnerisch und grafisch überprüfen.
  \end{itemize}
  \note{Person 1: Vorwissen abfragen: Was bedeutet eine Nullstelle?}
\end{frame}

\begin{frame}{Stundenplan und Aufgabenverteilung}
  % BEARBEITEN: Beispiel für 45 Minuten; an euren Zeitrahmen anpassen.
  \small
  \begin{tabular}{@{}p{1.4cm}p{4.2cm}p{5.3cm}@{}}
    \toprule
    Zeit & Phase & Zuständigkeit / Medium \\
    \midrule
    5 min & Einstieg und Vorwissen & Person 1 · Folie / Gespräch \\
    10 min & Gemeinsame Erarbeitung & Person 2 · Tafel \\
    15 min & Übung in Paaren & beide · Aufgabenfolie \\
    10 min & Ergebnisse sichern & Person 1 · Tafel, Person 2 moderiert \\
    5 min & Lernzielkontrolle & Person 2 · Abschlussfrage \\
    \bottomrule
  \end{tabular}
  \note{Die jeweils andere Person beobachtet Zeit und Fragen aus der Gruppe.}
\end{frame}

\section[Tafelphase]{Erarbeitung an der Tafel}
\begin{frame}{Mathematische Grundlage}
  % INLINE-MATHE: \( ... \); abgesetzte Formel: \[ ... \].
  Eine quadratische Funktion hat die Form
  \[
    f(x)=ax^2+bx+c, \qquad a,b,c\in\mathbb{R},\quad a\neq 0.
  \]
  \begin{block}{Definition: Nullstelle}
    Eine Zahl \(x_0\) heißt Nullstelle von \(f\), wenn \(f(x_0)=0\).
  \end{block}
  % QUELLEN: Schlüssel muss in literatur.bib existieren. Seiten ergänzen,
  % sobald ihr die verwendete Ausgabe und die betreffende Stelle geprüft habt.
  \small Hintergrund zur Analysis\footcite{spivak2008}.
  \note{Person 2: Definition an die Tafel schreiben. Beispiel selbst herleiten.}
\end{frame}

\begin{frame}{Tafelphase: gemeinsam herleiten}
  % BEARBEITEN: Diese Folie ist der Arbeitsauftrag. Erst rechnen lassen,
  % dann die nächste Folie zur Ergebnissicherung öffnen.
  \begin{block}{Arbeitsauftrag}
    Bestimmt die Nullstellen von \(f(x)=x^2-4x+3\).
  \end{block}
  \begin{enumerate}
    \item Was bedeutet \(f(x)=0\) für den Graphen?
    \item Welche zwei Zahlen haben Summe \(4\) und Produkt \(3\)?
    \item Schreibt \(f(x)\) als Produkt und nutzt die Nullproduktregel.
  \end{enumerate}
  \medskip
  \textbf{Tafelaufteilung:} links Begriff, Mitte Rechnung, rechts Skizze.
  \note{Person 2 schreibt, Person 1 sammelt Vorschläge. Die Rechnung zunächst
    ohne Projektionslösung entwickeln; Fehler gemeinsam besprechen.}
\end{frame}

\begin{frame}{Ergebnissicherung: Tafelbild in drei Schritten}
  % \pause blendet den nächsten Teil erst nach einem Klick ein.
  % Entfernen, wenn alles sofort sichtbar sein soll. Overlays sind keine Videos.
  \begin{columns}[T,onlytextwidth]
    \begin{column}{0.46\textwidth}
      \textbf{Rechnung}
      \begin{align*}
        0 &= x^2-4x+3 \\
          &= (x-1)(x-3)
      \end{align*}
      \pause
      Also: \(x_1=1\) oder \(x_2=3\).
      \medskip

      \textbf{Probe:} \(f(1)=0\), \(f(3)=0\).
    \end{column}
    \begin{column}{0.50\textwidth}
      % GRAFIK: Formel in addplot und Achsengrenzen passend ändern.
      \begin{tikzpicture}
        \begin{axis}[width=\linewidth,height=4.5cm,
          axis lines=middle,xlabel={$x$},ylabel={$f(x)$},
          xmin=0,xmax=4,ymin=-1.5,ymax=3.5,grid=major,
          xtick={0,1,2,3,4},ytick={-1,0,1,2,3}]
          \addplot[matheblau,thick,domain=0:4,samples=60]{x^2-4*x+3};
          \addplot[only marks,mark=*,mark size=2pt] coordinates {(1,0) (3,0)};
        \end{axis}
      \end{tikzpicture}
      \scriptsize Eigene Darstellung; Beispiel selbst berechnet.
    \end{column}
  \end{columns}
\end{frame}

\section[Übung]{Üben und sichern}
\begin{frame}{Partneraufgabe und kurze Lernzielkontrolle}
  % BEARBEITEN: Zeit, Schwierigkeit und Arbeitsauftrag an Lerngruppe anpassen.
  \begin{exampleblock}{Partnerarbeit · 5 Minuten}
    Bestimmt die Nullstellen von \(g(x)=x^2-5x+6\).
    Begründet jeden Schritt und prüft eure Lösung.
  \end{exampleblock}
  \begin{itemize}
    \item Person A rechnet, Person B erklärt und kontrolliert.
    \item Tauscht anschließend die Rollen.
  \end{itemize}
  \alert{Abschlussfrage:} Warum ist \(uv=0\) nur möglich, wenn mindestens
  einer der beiden Faktoren null ist?
  \note{Die Rollen A/B betreffen die Lernenden, nicht die Vortragenden.
    Musterlösung auf der nächsten Folie erst nach der Arbeitsphase zeigen.}
\end{frame}

\begin{frame}{Musterlösung und Merksatz}
  \begin{align*}
    g(x)&=x^2-5x+6=(x-2)(x-3), \\
    g(x)=0 &\iff x=2\ \text{oder}\ x=3.
  \end{align*}
  \begin{block}{Merksatz}
    Nullstellen sind die \(x\)-Koordinaten der Schnittpunkte des Graphen
    mit der \(x\)-Achse. Nicht jede quadratische Funktion hat reelle Nullstellen.
  \end{block}
  \note{Optionales Gegenbeispiel an der Tafel: h(x)=x^2+1.}
\end{frame}

\section{Quellen}
% ZITIERWERKSTATT: Alle folgenden Seitenangaben sind ungeprüfte Layoutbeispiele.
\begin{frame}{Zitierwerkstatt: Seiten direkt am Beleg}
  \begin{alertblock}{Nur eine Layoutdemonstration}
    Beispielsätze und Seitenzahlen sind keine geprüften Quellenbelege.
  \end{alertblock}
  \begin{itemize}
    \item Allgemeiner Verweis auf ein Buch\footcite{polya1945}.
    \item Eine konkrete Seite\footcite[S.~23]{polya1945}.
    \item Eine andere Stelle desselben Buches\footcite[S.~42]{polya1945}.
    \item Ein zusammenhängender Seitenbereich\footcite[S.~10--12]{spivak2008}.
  \end{itemize}
\end{frame}

\begin{frame}{Zitierwerkstatt: indirekt und direkt}
  \small\textbf{Layoutdemo:} Seiten ungeprüft; kein Originalzitat.
  \begin{itemize}
    \item Sinngemäße Beispielaussage mit Vorsatz\footcite[vgl.][S.~23--25]{polya1945}.
    \item Dasselbe mit einer zweiten Quelle\footcite[vgl.][S.~15]{spivak2008}.
    \item \enquote{Erfundener Demotext für die Darstellung eines direkten Zitats.}%
      \footcite[S.~23]{polya1945}
  \end{itemize}
  \medskip
  \small Echte direkte Zitate benötigen den geprüften Originalwortlaut
  und die tatsächliche Fundstelle.
\end{frame}

\begin{frame}{Zitierwerkstatt: mehrere Quellen und Spalten}
  \small\textbf{Layoutdemo:} Alle Seitenangaben sind ungeprüft.
  \medskip
  \begin{columns}[T,onlytextwidth]
    \begin{column}{0.47\textwidth}
      \textbf{Gemeinsamer Beleg}
      \par\medskip
      Eine Aussage kann zwei Bücher in einer Fußnote nennen%
      \footcite{polya1945,spivak2008}.
    \end{column}
    \begin{column}{0.47\textwidth}
      \textbf{Unterschiedliche Fundstellen}
      \par\medskip
      Jedes Buch bekommt seinen eigenen Seitenbeleg%
      \footcite[S.~23]{polya1945}\footcite[S.~10--12]{spivak2008}.
    \end{column}
  \end{columns}
  \medskip
  \small Die Quellen stehen auch bei Spalten unten auf der gesamten Folie.
\end{frame}

\begin{frame}[fragile]{Zitierwerkstatt: so setzt ihr die Fußmarken}
  \begin{verbatim}
Aussage\footcite{polya1945}.
Aussage\footcite[S.~23]{polya1945}.
Aussage\footcite[S.~23--25]{polya1945}.
Aussage\footcite[vgl.][S.~23--25]{polya1945}.
Aussage\footcite{polya1945,spivak2008}.
  \end{verbatim}
  \small Buchdaten in \texttt{literatur.bib}; Fundstellen direkt im Text.
  Die vollständigen Einträge folgen im Literaturverzeichnis.
\end{frame}

% LITERATUR: BibTeX erzeugt die Liste automatisch aus den zitierten Einträgen.
% Neue Quelle: Eintrag in literatur.bib + \footcite{schluessel} auf der Folie.
% allowframebreaks verteilt eine längere Liste automatisch auf weitere Folien.
\begin{frame}[allowframebreaks]{Literaturverzeichnis}
  \footnotesize
  \bibliographystyle{plainnat}
  \bibliography{literatur}
\end{frame}

% KOPIERVORLAGE für zusätzliche Folien (Prozentzeichen entfernen):
% \begin{frame}{Euer Folientitel}
%   Euer Inhalt, z. B. \(a^2+b^2=c^2\).
%   \note{Person 1/2: Hinweis nur für euch.}
% \end{frame}
% BILDDATEI: Bild im Ordner ablegen, dann z. B.:
% \includegraphics[width=0.7\textwidth]{bild.pdf}
% Direkt daneben Urheber, Quelle (\footcite{...}) und ggf. Lizenz nennen.
\end{document}
